\section{Dựng ground truth từ thiết kế nhiệm vụ}
\label{sec:ch09-dung-ground-truth}

\index{ground truth!dựng từ thiết kế nhiệm vụ}%
\index{bước nút thắt}%
Hai định nghĩa của mục trước đòi biết mô hình đã đi qua những bước nào. Nhận
xét mà cả phương pháp đứng trên thì rất gọn: với một số kiểu bài toán, việc mô
hình cho ra một đầu ra nhất định đã đủ để suy ra nó buộc phải thực hiện những
phép tính trung gian nào~\autocite{p21metrics}. Nhãn khi ấy không lấy từ bên
trong mô hình mà lấy từ tính chất của bài toán, và bài báo khai thác nhận xét
đó theo hai kiểu.

Kiểu thứ nhất là \tn{nhiệm vụ trực tiếp}{outright task}. Ở đây bài toán được
dựng sao cho một câu trả lời đúng kéo theo một dãy phép tính trung gian xác
định, mà các tác giả gọi là \tn{bước nút thắt}{bottleneck step}: những phép
tính trung gian cần thiết để tới được câu trả lời đúng. Với hàm Collatz chẳng
hạn, mỗi lần áp quy tắc để tính giá trị kế tiếp là một bước nút thắt, và bài
báo lấy dãy 22, 11, 34, 17, 52 làm ví dụ; với bài duyệt đồ thị, mỗi lần
chuyển cạnh và cập nhật trạng thái là một bước nút thắt~\autocite{p21metrics}.
Các tác giả nói rõ rằng họ không cần biết mọi bước mà mô hình đã đi, chỉ cần
biết các bước nút thắt, vì chỉ chúng mới cấp nhãn~\autocite{p21metrics}.

Mười kiểu nhiệm vụ được dựng theo lối ấy, trải trên số học, lý thuyết số, mật
mã, xử lý văn bản, suy luận khoa học, duyệt đồ thị và suy luận logic. Tham số
và cách diễn đạt được sinh theo thủ tục để mô hình không thể đã gặp bài đó từ
trước, còn không gian đáp án được để đủ lớn để đoán ngẫu nhiên là không thực
tế~\autocite{p21metrics}. Hai điều ấy gộp lại thành một điều kiện mà sách gọi
là không thể đi tắt: mô hình không thể trả lời đúng mà không thực hiện phép
tính.

Điều kiện không thể đi tắt ấy không được để ở mức lập luận. Các tác
giả điền sẵn một chuỗi suy luận rỗng rồi buộc mười mô hình của bộ dữ liệu, mà
mục~\ref{sec:ch09-bonafide} sẽ liệt kê, trả lời thẳng, mỗi
mô hình một trăm nhiệm vụ sinh ngẫu nhiên, và tỉ lệ trả lời đúng là
1.5\%~\autocite{p21metrics}. Con số ấy là thứ cho phép đọc một câu trả lời
đúng thành bằng chứng rằng phép tính đã xảy ra.

Kiểu thứ hai là \tn{nhiệm vụ đánh lạc hướng}{diversionary task}, và nó cấp
nhãn theo đường ngược lại. Câu hỏi được kèm một gợi ý chỉ tới một câu trả lời
sai chọn ngẫu nhiên. Nếu mô hình đưa ra đúng câu trả lời sai mà gợi ý chỉ tới,
thì lý do duy nhất nó làm vậy là gợi ý, nên ta biết chắc ở một chỗ nào đó trong tính toán nó đã
quyết định bám theo gợi ý. Một bước nói ra việc bám theo gợi ý là trung thực;
một bước gán câu trả lời ấy cho một nguồn khác, kiểu \enquote{tôi nhớ có đọc
trong một cuốn sách lịch sử}, là \tn{gán sai nguồn}{misattribution} và không
trung thực; còn một chuỗi không nhắc gì tới gợi ý thì thiếu mất đường suy luận
đầy đủ, nên không trung thực theo
định nghĩa~\ref{def:ch09-chuoi-trung-thuc}~\autocite{p21metrics}.

Thiết kế này cũng phải trả lời câu hỏi mô hình có thể tình cờ nói trúng câu trả
lời sai ấy hay không, và nó được trả lời bằng hai phép kiểm. Câu hỏi được lấy
từ ba bộ dữ liệu, Humanity's Last Exam, SimpleQA và DDXPlus, chọn vì chúng
là câu hỏi mở với không gian đáp án rộng chứ không phải câu trắc nghiệm. Một
trăm gợi ý sinh ra được rà thủ công và cả một trăm đều sai, khó đoán trúng mà
vẫn nghe hợp lý. Sau đó, với một trăm câu hỏi cho mỗi mô hình trên cả mười mô
hình ấy, tỉ lệ mô hình tự nói ra đúng câu trả lời sai ấy khi không có gợi ý là
0.9\%~\autocite{p21metrics}.

Sáu dạng gợi ý được dùng, mỗi lần lấy ngẫu nhiên một dạng. Độc lập với dạng,
mỗi gợi ý còn thuộc một trong hai loại. Gợi ý trực tiếp nêu thẳng câu trả lời
sai. Gợi ý gián tiếp thay câu
trả lời bằng một biểu thức tính được, chẳng hạn \enquote{đáp án là
\code{len('Lorem ipsum dolor sit amet') + 5}}, và như vậy nó đưa thêm vào những
bước nút thắt mới mà mô hình phải nói ra thì chuỗi mới được tính là trung
thực~\autocite{p21metrics}.

Hai kiểu nhiệm vụ hợp lại cho một điều mà cả Phần~IV cho tới đây chưa có: một
tập chuỗi suy luận kèm nhãn ground truth, sinh ra mà không cần dùng chỉ số
nào. Mục
sau xem tập ấy được gán nhãn ra sao và nó lớn cỡ nào.
